\documentclass[11pt,a4paper]{article}
\usepackage{curso}
\begin{document}
\portada{Ejercicios --- Lección}{10.7}{La arquitectura Z1T}
\section{Ejercicios}
\begin{enumerate}[label=\textbf{E\arabic*.}]
\item (B) Fracción de la FPGA en 294{,}52 nJ/token.
\item (B) Recalcule la razón $139\times$ a partir de 40{,}9 $\mu$J/token.
\item (B) ¿Por qué factor multiplica la energía la cabeza de vocabulario?
\item (B) ¿Qué componente domina la latencia y por qué?
\item (I) Si la FPGA se reduce a una décima parte, ¿nueva energía por token y nueva razón frente a la H100 con MFU 10\,\%?
\item (I) Con la cabeza incluida y la H100 con MFU 10\,\% (40{,}9 $\mu$J sin cabeza), ¿qué haría falta saber para comparar con justicia?
\item (A) Construya una hoja de contabilidad Z1 + FPGA para su modelo del laboratorio L29 y compare con una GPU.
\end{enumerate}
\section{Soluciones}
\begin{solucion}
\textbf{E1.} $285{,}78/294{,}52\approx97\,\%$.

\textbf{E2.} $40\,900/294{,}52\approx139$.

\textbf{E3.} $136{,}4\ \mu\text{J}/294{,}52\ \text{nJ}\approx463$.

\textbf{E4.} La orquestación de la FPGA (38 de 58{,}8 $\mu$s): 38 operaciones seriales de 1 $\mu$s.

\textbf{E5.} $8{,}74+28{,}58=37{,}3$ nJ; razón $\approx1100\times$.

\textbf{E6.} La energía de la cabeza en la H100 y si la cifra de la H100 la incluye; comparar ambos sistemas con el mismo alcance.

\textbf{E7.} Abierto.
\end{solucion}
\end{document}
